\documentclass[12pt,a4paper]{ctexart}

\usepackage[a4paper,left=2.6cm,right=2.6cm,top=1.55cm,bottom=1.6cm]{geometry}
\usepackage{xcolor}
\usepackage{fancyhdr}
\usepackage{setspace}
\usepackage{parskip}
\usepackage{hyperref}

\definecolor{mainblue}{RGB}{32,83,154}
\definecolor{lightblue}{RGB}{236,243,252}
\definecolor{softgray}{RGB}{245,246,248}

\hypersetup{colorlinks=false,pdfborder={0 0 0}}
\pagestyle{fancy}
\fancyhf{}
\fancyhead[L]{\small Presentation\_Group8 逐字稿}
\fancyhead[R]{\small 编译原理大作业}
\fancyfoot[C]{\small \thepage}
\renewcommand{\headrulewidth}{0.4pt}
\renewcommand{\footrulewidth}{0pt}
\setlength{\headheight}{15pt}
\setstretch{1.42}
\setlength{\parindent}{0pt}
\setlength{\parskip}{0.55em}

\newcommand{\code}[1]{\texttt{\detokenize{#1}}}
\newcommand{\speaker}[1]{%
  \vspace{0.65em}
  \noindent\colorbox{mainblue}{\parbox{\dimexpr\linewidth-2\fboxsep\relax}{\color{white}\bfseries #1}}
  \vspace{0.25em}
}
\newcommand{\cue}[1]{%
  \par\smallskip
  \noindent\colorbox{lightblue}{\parbox{\dimexpr\linewidth-2\fboxsep\relax}{\color{mainblue}\bfseries #1}}
  \par\smallskip
}
\newcommand{\qa}[1]{%
  \par\smallskip
  \noindent\colorbox{softgray}{\parbox{\dimexpr\linewidth-2\fboxsep\relax}{\bfseries #1}}
  \par\smallskip
}

\begin{document}

\begin{center}
{\LARGE\bfseries Presentation\_Group8 逐字稿}\\[-0.1em]
{\large 类 Rust 编译前端与四元式中间代码生成}\\[0.35em]
主讲顺序：冉子易 $\rightarrow$ 高胜寒 $\rightarrow$ 李堂辉 \quad
主讲约 6 分 50 秒，预留约 1 分钟问答
\end{center}

\speaker{冉子易：开场与大作业 1，第 1--8 页，约 2 分 25 秒}

\cue{第 1 页，封面，约 10 秒}
各位老师好，我们是第 8 组。汇报题目是“编译原理大作业 1 和 2”，内容是类 Rust 语言的前端分析和四元式中间代码生成。封面姓名按学号顺序排列。

\cue{翻到第 2 页，主线，约 20 秒}
对标老师的两个要求 PPTX，这两份作业可以串成一条前端流程。大作业 1 负责 Token 和 AST；大作业 2 继续消费 AST，做语义检查，并生成四元式 IR。

\cue{翻到第 3 页，对标要求，约 20 秒}
这一页对应老师两个要求 PPTX。必做语法规则覆盖到 5.1；扩展结构包括引用、数组、元组、表达式块和三类循环；大作业 2 重点完成语义诊断和四元式 IR 生成。

\cue{翻到第 4 页，AST 接口，约 20 秒}
这一页是两份作业衔接的核心。作业 1 输出的 AST 保留了类型、左值、控制流和表达式形态。作业 2 直接按这些节点做类型检查、循环检查，并继续生成 IR。

\cue{翻到第 5 页，扩展语法结构，约 20 秒}
老师 PPTX 中 0.1 到 5.1 是最低要求。我们比较值得讲的是后续扩展规则：引用、可变引用、数组、元组、表达式块、\code{if} 表达式和三类循环，都被保留成后续可以分析的结构。

\cue{翻到第 6 页，Parser 层次，约 20 秒}
Parser 使用手写递归下降。顶层解析函数、代码块和语句；表达式层再处理区间、比较、加减、乘除、调用、索引和引用。这样的层次方便维护，也方便定位错误。

\cue{翻到第 7 页，AST 设计，约 20 秒}
AST 保留后续阶段需要的信息。\code{Type} 里有 \code{I32}、引用、可变引用、数组和元组；\code{Expr} 里有二元运算、引用、索引、表达式块和 \code{if} 表达式。

\cue{翻到第 8 页，作业 1 产出，约 20 秒}
作业 1 的产出是可继续编译的结构。Token 提供类别和位置，AST 提供语句、表达式和类型结构。下面由高胜寒介绍语义分析。

\speaker{高胜寒：语义分析，第 9--12 页，约 1 分 30 秒}

\cue{翻到第 9 页，语义状态，约 20 秒}
第二份作业的特色是带状态遍历 AST。\code{Compiler} 维护 \code{scopes}、\code{current_return_type}、\code{loop_stack}，以及临时变量和标签计数器。语义是否合法，IR 怎么落地，都由这些状态一起决定。

\cue{翻到第 10 页，符号表栈，约 20 秒}
符号表中的变量记录可变性、静态类型和 IR 参数。作用域用 \code{Vec<HashMap<...>>} 表示，进入块时压栈，退出块时弹栈，查找变量时从内层向外层查。

\cue{翻到第 11 页，检查点，约 25 秒}
这一页对应老师要求 PPTX 里容易被追问的语义点：类型要对应，不可变变量不能二次赋值，\code{break}、\code{continue} 必须在循环内。实现上分别落到 \code{Type} 枚举、符号表里的 \code{mutable} 字段和 \code{loop_stack}。

\cue{翻到第 12 页，错误汇总，约 25 秒}
错误会集中收集到 \code{errors: Vec<String>}，最后统一输出。因此同一段程序可以同时报告不可变变量赋值、循环外 \code{break} 和返回类型不一致。下面由李堂辉介绍 IR 生成。

\speaker{李堂辉：IR 生成、验证与收束，第 13--19 页，约 2 分 45 秒}

\cue{翻到第 13 页，IR 结构，约 20 秒}
大作业 2 要求实现中间代码生成，并建议用四元式。我们的 IR 每条指令包含 \code{op}、\code{arg1}、\code{arg2} 和 \code{result}；\code{Arg} 可以是空值、整数、变量、临时变量或标签。

\cue{翻到第 14 页，表达式翻译，约 25 秒}
表达式翻译时，\code{compile_expression} 返回 \code{(Type, Arg)}。例如 \code{a + b * 2} 先生成 \code{b * 2} 的 \code{t0}，再生成 \code{a + t0} 的 \code{t1}，最后把 \code{t1} 赋值给 \code{c}。

\cue{翻到第 15 页，分支翻译，约 20 秒}
分支结构用条件跳转和标签表示。先计算条件，再用 \code{JUMP_IF_FALSE} 跳到 else 标签；then 分支结束后跳到结束标签。

\cue{翻到第 16 页，循环翻译，约 20 秒}
循环结构依赖入口标签、出口标签和循环栈。\code{continue} 跳入口，\code{break} 跳出口。编译循环体时把两个标签压入 \code{loop_stack}。

\cue{翻到第 17 页，扩展结构落点，约 30 秒}
这一页讲扩展结构在第二份作业里的真实落点。引用生成 \code{REF}，解引用生成 \code{DEREF}，数组读取生成 \code{LOAD_ARRAY}。元组下标会根据元组类型推导元素类型；函数调用先生成参数 \code{ARG}，再生成 \code{CALL}。函数签名检查和数组写入目前作为后续扩展说明。

\cue{翻到第 18 页，测试覆盖，约 25 秒}
测试也按老师规则组织。\code{test_all.rs.txt} 覆盖主要语法结构，\code{test_errors.rs.txt} 覆盖词法和语法错误，\code{test_semantic_errors.rs.txt} 覆盖未声明、类型不匹配、不可变赋值和非法控制流等语义错误。

\cue{翻到第 19 页，总结，约 30 秒}
总结一下，两份作业完成了从类 Rust 源程序到四元式 IR 的前端流程。作业 1 把语法结构化，作业 2 带状态遍历 AST，检查作用域、类型、可变性、返回类型和循环上下文。最终 IR 用临时变量和标签线性化表达式和控制流。我们的汇报到这里结束，谢谢老师，欢迎提问。

\clearpage
\speaker{备用问答提示，不计入主讲时间}

\qa{如果老师问两份作业区别}
大作业 1 的输入是源代码字符串，核心状态是 Token 指针和解析函数栈，输出是 Token 和 AST。大作业 2 的输入是 AST，核心状态是符号表栈、循环栈、返回类型和计数器，输出是语义错误或四元式 IR。

\qa{如果老师问实现边界}
当前重点是前端流程和四元式表达，没有继续做目标代码生成。函数调用已经生成 \code{ARG/CALL}，但函数签名检查还可以继续完善。数组读取有 \code{LOAD_ARRAY}，数组写入接口存在，后续可以扩展更完整的地址计算和 \code{STORE_ARRAY}。

\qa{如果老师问实现特色}
第一，两份作业用 AST 接成一条流程；第二，引用、数组、元组、表达式块和循环结构都进入 AST；第三，语义分析带着 \code{scopes}、\code{current_return_type} 和 \code{loop_stack} 运行；第四，IR 用临时变量和标签把表达式、分支和循环线性化。这些点都能在 \code{ast.rs}、\code{compiler.rs} 和 \code{ir.rs} 中对应到代码。

\qa{如果老师问和要求 PPTX 的对应关系}
必做语法规则覆盖到老师要求的 5.1；扩展规则中，引用、数组、元组、表达式块、循环控制进入 AST，并在语义或 IR 阶段有对应处理。大作业 2 对应语义诊断和四元式生成，重点实现了类型一致、可变性、返回类型、循环控制和标签化 IR。

\end{document}
